% !TEX program = lualatex
% 한국교원대학교 학위논문 Beamer 발표 템플릿의 메인 파일입니다.
% Main file for the KNUE thesis/dissertation Beamer presentation template.
%
% 이 파일은 발표 전체의 순서와 공통 설정을 조립합니다.
% This file assembles the presentation order and shared settings.
%
% 실제 발표 내용은 sub/ 폴더의 장별 파일에서 수정합니다.
% Edit the actual presentation content in the section files under sub/.
%
% 빌드 명령:
% Build command:
%   latexmk -pdf -lualatex KNUE_beamer_main.tex

% hyperref가 Beamer 발표자 노트나 중복 페이지에서 같은 링크 이름을 만들지 않게 합니다.
% Prevent hyperref from creating duplicate hyperlink destinations in Beamer notes or repeated pages.
\PassOptionsToPackage{hypertexnames=false}{hyperref}

% 왼쪽 사이드바 목차에는 현재 섹션의 하위 항목만 자세히 보이도록 설정합니다.
% Show detailed sidebar subsections only for the current section.
\PassOptionsToPackage{hideothersubsections}{beamerouterthemesidebar}

% KNUE 발표용으로 만든 Beamer 클래스 파일을 불러옵니다.
% Load the custom Beamer class for KNUE presentations.
\documentclass{KNUE_beamer_class}

% 그림 검색 경로입니다. figures/images는 images/에, 로고는 logo/에 넣는 것을 권장합니다.
% Graphic search paths. Put figures/images in images/ and logos in logo/.
\graphicspath{{images/}{logo/}}

% 참고문헌 파일 경로입니다. 한국어/영어 문헌을 sub/references.bib 하나에서 관리합니다.
% Bibliography file path. Korean and English references are managed together in sub/references.bib.
\providecommand{\beamerbib}{sub/references.bib}

% 제목, 저자, 참고문헌, 발표자 노트, 편의 명령을 모아 둔 공통 설정 파일입니다.
% Shared settings for title metadata, bibliography, speaker notes, and helper commands.
% 한국교원대학교 Beamer 발표 템플릿의 공통 설정 파일입니다.
% Shared settings for the KNUE Beamer presentation template.
%
% 제목, 저자, 참고문헌, 발표자 노트, 자주 쓰는 명령을 이 파일에서 관리합니다.
% This file manages title metadata, bibliography, speaker notes, and reusable helper commands.

% ========== Presentation Information / 발표 정보 ==========

% 아래 명령의 중괄호 안만 자신의 발표 정보로 바꿉니다.
% Replace only the text inside braces with your own presentation information.

% 학위과정 구분입니다. 논문 템플릿처럼 석사 또는 박사 중 하나로 바꿉니다.
% Degree program selector. As in the thesis template, choose either 석사 or 박사.
\providecommand{\degreeProgramKor}{박사} % 석사 또는 박사

% 발표 종류 구분입니다. proposal은 논문 프로포절, result는 논문 결과발표입니다.
% Presentation-stage selector. Use proposal for proposal talks and result for result/defense talks.
\providecommand{\presentationStage}{proposal} % proposal 또는 result

% 선택값을 preamble에서 한 번만 판정합니다. 한글 선택값은 토큰을 문자열로 바꿔 비교해야 안정적입니다.
% Evaluate selector values once in the preamble. Korean selector values are compared as strings for stability.
\usepackage{pdftexcmds}
\makeatletter
\newif\ifKNUEMaster
\newif\ifKNUEResult
\edef\KNUE@mastervalue{\detokenize{석사}}
\edef\KNUE@resultvalue{\detokenize{result}}
\edef\KNUE@degreevalue{\detokenize\expandafter{\degreeProgramKor}}
\edef\KNUE@stagevalue{\detokenize\expandafter{\presentationStage}}
\ifnum\pdf@strcmp{\KNUE@degreevalue}{\KNUE@mastervalue}=0
  \KNUEMastertrue
\else
  \KNUEMasterfalse
\fi
\ifnum\pdf@strcmp{\KNUE@stagevalue}{\KNUE@resultvalue}=0
  \KNUEResulttrue
\else
  \KNUEResultfalse
\fi
\makeatother

% 학위과정에 따라 국문 학위논문 문구를 만듭니다.
% Build the Korean thesis/dissertation label from the degree program.
\newcommand{\BeamerDegreeWorkKor}{\degreeProgramKor{}학위논문}

% 학위과정에 따라 영문 학위논문 문구를 미리 정합니다.
% Predefine the English thesis/dissertation label from the degree program.
\ifKNUEMaster
  \newcommand{\BeamerDegreeWorkEng}{Master's Thesis}%
\else
  \newcommand{\BeamerDegreeWorkEng}{Doctoral Dissertation}%
\fi

% 발표 종류에 따라 국문/영문 발표 문구를 미리 정합니다.
% Predefine Korean/English presentation-stage labels from the stage selector.
\ifKNUEResult
  \newcommand{\BeamerStageKor}{결과발표}%
  \newcommand{\BeamerStageEng}{Results Presentation}%
\else
  \newcommand{\BeamerStageKor}{연구계획 발표}%
  \newcommand{\BeamerStageEng}{Proposal Presentation}%
\fi

% 발표 종류에 따라 두 가지 내용 중 하나를 골라 출력합니다.
% Choose between two texts according to the selected presentation stage.
\newcommand{\KNUEStageText}[2]{%
  \ifKNUEResult #2\else #1\fi%
}

% 발표 종류에 따라 결과 섹션의 제목 문구를 미리 정합니다.
% Predefine result-section labels according to the selected presentation stage.
\ifKNUEResult
  \newcommand{\BeamerResultsSectionTitle}{연구 결과}%
  \newcommand{\BeamerResultsSectionTitleEng}{Research Results}%
  \newcommand{\BeamerResultsFirstSubsection}{주요 결과}%
  \newcommand{\BeamerResultsFirstFrame}{주요 연구 결과}%
  \newcommand{\BeamerResultsFirstBlock}{연구를 통해 확인한 결과}%
  \newcommand{\BeamerResultsSecondSubsection}{결과 해석}%
  \newcommand{\BeamerResultsSecondFrame}{결과 해석}%
  \newcommand{\BeamerResultsThirdSubsection}{논의와 시사점}%
  \newcommand{\BeamerResultsThirdFrame}{논의와 시사점}%
  \newcommand{\BeamerSummaryFrame}{연구 결과 요약}%
  \newcommand{\BeamerQASubsection}{질의 및 논의}%
  \newcommand{\BeamerQAFrame}{질의 및 논의 사항}%
  \newcommand{\BeamerQABlock}{심사위원과 논의할 부분}%
  \newcommand{\BeamerDefaultShortTitle}{Results}%
\else
  \newcommand{\BeamerResultsSectionTitle}{예상 결과}%
  \newcommand{\BeamerResultsSectionTitleEng}{Expected Results}%
  \newcommand{\BeamerResultsFirstSubsection}{예상 산출물}%
  \newcommand{\BeamerResultsFirstFrame}{예상 산출물}%
  \newcommand{\BeamerResultsFirstBlock}{연구를 통해 산출할 결과}%
  \newcommand{\BeamerResultsSecondSubsection}{연구 일정}%
  \newcommand{\BeamerResultsSecondFrame}{연구 일정}%
  \newcommand{\BeamerResultsThirdSubsection}{기대 효과}%
  \newcommand{\BeamerResultsThirdFrame}{기대 효과}%
  \newcommand{\BeamerSummaryFrame}{연구계획 요약}%
  \newcommand{\BeamerQASubsection}{검토 요청}%
  \newcommand{\BeamerQAFrame}{검토 요청 사항}%
  \newcommand{\BeamerQABlock}{심사위원께 검토를 요청할 부분}%
  \newcommand{\BeamerDefaultShortTitle}{Proposal}%
\fi

% 국문 발표 제목입니다. 표지에서 가장 크게 출력됩니다.
% Korean presentation title. This is shown most prominently on the title slide.
\providecommand{\BeamerTitleKor}{학위논문 발표 제목}

% 영문 발표 제목입니다. 국문 제목 아래에 작은 글씨로 출력됩니다.
% English presentation title. This appears below the Korean title in smaller type.
\providecommand{\BeamerTitleEng}{Thesis or Dissertation Presentation Title}

% 발표의 성격을 짧게 설명하는 국문 부제입니다.
% Korean subtitle that describes the presentation type.
\providecommand{\BeamerSubtitleKor}{\BeamerDegreeWorkKor{} \BeamerStageKor}

% 발표의 성격을 짧게 설명하는 영문 부제입니다.
% English subtitle that describes the presentation type.
\providecommand{\BeamerSubtitleEng}{\BeamerDegreeWorkEng{} \BeamerStageEng}

% PDF 북마크와 사이드바에서 사용할 짧은 제목입니다.
% Short title used in PDF bookmarks and the sidebar.
\providecommand{\BeamerShortTitle}{\BeamerDefaultShortTitle}

% 발표자 이름입니다.
% Presenter name.
\providecommand{\BeamerAuthor}{홍 길 동}

% 지도교수 이름입니다. 필요하면 "지도교수:" 문구까지 함께 수정합니다.
% Advisor name. Edit the "지도교수:" label too if needed.
\providecommand{\BeamerAdvisor}{지도교수: 성 춘 향}

% 소속 학과, 전공, 연구실 정보를 적습니다.
% Department, major, lab, or institutional affiliation.
\providecommand{\BeamerInstitute}{한국교원대학교 지구과학교육전공}

% 발표일입니다. 기본값 \today는 빌드한 날짜를 자동으로 넣습니다.
% Presentation date. The default \today inserts the build date automatically.
\providecommand{\BeamerDate}{\today}

% 이전 연구계획서 템플릿 명령을 계속 사용할 수 있게 남겨 둔 호환 별칭입니다.
% Compatibility aliases for documents created from the earlier proposal template.
\providecommand{\ProposalTitleKor}{\BeamerTitleKor}
\providecommand{\ProposalTitleEng}{\BeamerTitleEng}
\providecommand{\ProposalSubtitle}{\BeamerSubtitleKor}
\providecommand{\ProposalAuthor}{\BeamerAuthor}
\providecommand{\ProposalAdvisor}{\BeamerAdvisor}
\providecommand{\ProposalInstitute}{\BeamerInstitute}
\providecommand{\ProposalDate}{\BeamerDate}

% ========== Packages / 패키지 ==========

% multicol은 긴 목차를 두 단으로 나눌 때 사용합니다.
% multicol is used to split long table-of-contents slides into two columns.
\usepackage{multicol}

% csquotes는 biblatex-apa와 함께 인용부호 처리를 안정적으로 해 줍니다.
% csquotes works with biblatex-apa to handle quotation marks consistently.
\usepackage{csquotes}

% APA 7판 형식 참고문헌을 biber로 처리합니다.
% Use biber to process APA 7th edition references.
\usepackage[
  backend=biber,
  style=apa,
  natbib=true,
  maxbibnames=20,
  minbibnames=20,
  maxcitenames=2,
  mincitenames=1,
  uniquelist=false,
  dashed=false,
  date=year,
  urldate=iso,
  dateera=astronomical,
  datezeros=true,
  seconds=true
]{biblatex}

% 논문 템플릿과 동일하게 biblatex-apa의 APA 7판 지역화 규칙을 사용합니다.
% Use the same APA 7 localization rules as the thesis template.
\DeclareLanguageMapping{english}{english-apa}
\DeclareLanguageMapping{korean}{english-apa}

% 참고문헌 파일 경로입니다. 메인 파일에서 다른 경로를 지정하면 그 값을 우선 사용합니다.
% Bibliography path. If the main file defines another path, that value takes precedence.
\providecommand{\beamerbib}{sub/references.bib}
\providecommand{\proposalbib}{\beamerbib}

% 실제 참고문헌 데이터베이스를 biblatex에 등록합니다.
% Register the bibliography databases with biblatex.
\addbibresource{\beamerbib}

% pgfpages는 발표자 노트를 두 번째 화면에 배치할 때 필요합니다.
% pgfpages is needed when speaker notes are placed on a second screen.
\usepackage{pgfpages}

% ========== Bibliography and Citation Settings / 참고문헌과 인용 설정 ==========

% references.bib의 langid=korean 문헌을 먼저, langid=english 문헌을 나중에 정렬합니다.
% Sort langid=korean entries first and langid=english entries later.
\DeclareSourcemap{
  \maps[datatype=bibtex]{
    \map{
      \step[fieldsource=langid, match=\regexp{korean}, final]
      \step[fieldset=presort, fieldvalue={a}]
    }
    \map{
      \step[fieldsource=langid, match=\regexp{english}, final]
      \step[fieldset=presort, fieldvalue={z}]
    }
  }
}

% 참고문헌 목록에서 한국어 문헌은 저자명을 가운데점으로 연결하고, 영어 문헌은 APA 7판 기본 연결 방식을 사용합니다.
% Korean bibliography entries join author names with centered dots; English entries keep APA 7 delimiters.
\DeclareDelimFormat[bib,biblist]{multinamedelim}{%
  \iffieldequalstr{langid}{korean}
  {\textperiodcentered}
  {\addcomma\space}%
}
\DeclareDelimFormat[bib,biblist]{finalnamedelim}{%
  \iffieldequalstr{langid}{korean}
  {\textperiodcentered}
  {\ifnumgreater{\value{liststop}}{2}{\finalandcomma}{}\addspace\&\space}%
}
\DeclareDelimFormat[bib,biblist]{nameyeardelim}{%
  \iffieldequalstr{langid}{korean}
  {}
  {\newunitpunct}%
}
\DeclareDelimAlias[bib,biblist]{nonameyeardelim}{nameyeardelim}

% 한국어 학술지명과 권은 굵게, 영어 학술지명과 권은 APA 기본 기울임꼴로 표시합니다.
% Korean journal titles and volumes are bold; English ones remain APA-style italics.
\DeclareFieldFormat[article]{journaltitle}{%
  \iffieldequalstr{langid}{korean}
  {\mkbibbold{#1}}
  {\mkbibemph{#1}}%
}
\DeclareFieldFormat[article]{volume}{%
  \iffieldequalstr{langid}{korean}
  {\mkbibbold{#1}}
  {\mkbibemph{#1}}%
}

% DOI와 URL은 클릭 가능한 주소 형태로 표시합니다.
% Display DOI and URL fields as clickable address-style links.
\DeclareFieldFormat{doi}{%
  \ifhyperref{\href{https://doi.org/#1}{\nolinkurl{https://doi.org/#1}}}{\nolinkurl{https://doi.org/#1}}%
  \nopunct%
}
\DeclareFieldFormat{url}{\url{#1}}

% 서술형 인용에서 저자명과 연도 사이 공백을 줄입니다.
% Remove the space between author names and year in textual citations.
\DeclareDelimFormat[textcite]{nameyeardelim}{}

% Beamer 참고문헌은 본문보다 작지만 읽을 수 있도록 scriptsize로 설정합니다.
% Use scriptsize for Beamer bibliography entries so they remain readable.
\renewcommand*{\bibfont}{\scriptsize}

% 참고문헌의 내어쓰기 폭을 줄여 좁은 슬라이드에서도 줄바꿈이 덜 부담스럽게 보이게 합니다.
% Use a shorter hanging indent so wrapped bibliography lines fit Beamer slides better.
\setlength{\bibhang}{0.8em}

% 참고문헌 항목 사이의 간격을 줄여 한 슬라이드에 더 많이 들어가게 합니다.
% Reduce spacing between bibliography entries so more items fit on slides.
\setlength\bibitemsep{0pt}
\setlength\bibnamesep{0pt}

% 참고문헌 아이콘 대신 텍스트형 항목 표시를 사용합니다.
% Use text-style bibliography items instead of icon bullets.
\setbeamertemplate{bibliography item}[text]

% 참고문헌 제목을 한국어로 보이게 합니다.
% Display bibliography headings in Korean.
\DefineBibliographyStrings{english}{
  references = {참고문헌},
  bibliography = {참고문헌},
}

% ---------- 한국어 문헌 본문 인용 명령 / Korean citation commands ----------
% 한국어 문헌은 \ktextcite, \kparencite를 사용하고 영어 문헌은 \textcite, \parencite를 사용합니다.
% Use \ktextcite and \kparencite for Korean sources; use \textcite and \parencite for English sources.
% \kparencite와 \parencite는 괄호를 직접 출력하지 않습니다.
% \kparencite and \parencite do not print outer parentheses; write them in the slide text when needed.

\newcommand{\KoreanCiteAuthorDelim}{\textperiodcentered}
\newcommand{\UseKoreanCiteDelims}{%
  \DeclareDelimFormat{multinamedelim}{\KoreanCiteAuthorDelim}%
  \DeclareDelimFormat[textcite]{multinamedelim}{\KoreanCiteAuthorDelim}%
  \DeclareDelimFormat[parencite]{multinamedelim}{\KoreanCiteAuthorDelim}%
  \DeclareDelimFormat{finalnamedelim}{\KoreanCiteAuthorDelim}%
  \DeclareDelimFormat[textcite]{finalnamedelim}{\KoreanCiteAuthorDelim}%
  \DeclareDelimFormat[parencite]{finalnamedelim}{\KoreanCiteAuthorDelim}%
  \DeclareDelimFormat[textcite]{nameyeardelim}{}%
  \DeclareDelimFormat{andothersdelim}{\addspace}%
  \DeclareDelimFormat[textcite]{andothersdelim}{\addspace}%
  \DeclareDelimFormat[parencite]{andothersdelim}{\addspace}%
  \csdef{abx@lstr@andothers}{외}%
  \csdef{abx@sstr@andothers}{외}%
}
\let\KNUEorigparencite\parencite
\let\parencite\cite
\NewDocumentCommand{\ktextcite}{o o m}{%
  \begingroup
  \UseKoreanCiteDelims
  \IfNoValueTF{#1}{%
    \textcite{#3}%
  }{%
    \IfNoValueTF{#2}{\textcite[#1]{#3}}{\textcite[#1][#2]{#3}}%
  }%
  \endgroup
}
\NewDocumentCommand{\kparencite}{o o m}{%
  \begingroup
  \UseKoreanCiteDelims
  \IfNoValueTF{#1}{%
    \cite{#3}%
  }{%
    \IfNoValueTF{#2}{\cite[#1]{#3}}{\cite[#1][#2]{#3}}%
  }%
  \endgroup
}

% ========== Table of Contents / 목차 설정 ==========

% 목차에서 section과 subsection 앞에 붙일 점 표식입니다.
% Bullet marker used before section and subsection entries in the table of contents.
\newcommand{\knueTocBullet}{\raisebox{0.15ex}{\large\textbullet}}

% 목차에서 section은 subsection보다 한 단계 크게 표시합니다.
% Sections are displayed one size larger than subsections in ToC slides.
\setbeamerfont{section in toc}{size=\small}
\setbeamerfont{subsection in toc}{size=\scriptsize}

% section 목차 항목의 줄바꿈과 들여쓰기를 안정화합니다.
% Stabilize wrapping and indentation for section entries in the ToC.
\setbeamertemplate{section in toc}{%
  \vspace{0.35em}%
  \usebeamerfont{section in toc}%
  \usebeamercolor[fg]{section in toc}%
  \leavevmode\parindent0pt%
  \hangindent=1.4em\hangafter=1%
  \makebox[1.2em][l]{\knueTocBullet}\inserttocsection\par%
}

% subsection 목차 항목은 section보다 안쪽에 들여쓰기합니다.
% Subsection entries are indented relative to section entries.
\setbeamertemplate{subsection in toc}{%
  \vspace{0.15em}%
  \usebeamerfont{subsection in toc}%
  \usebeamercolor[fg]{subsection in toc}%
  \leavevmode\parindent0pt%
  \leftskip=1.4em%
  \hangindent=1.4em\hangafter=1%
  \makebox[1.2em][l]{\knueTocBullet}\inserttocsubsection\par%
}

% ========== Presenter Notes / 발표자 노트 ==========

% overlay에서 아직 드러나지 않은 항목을 희미하게 보이도록 설정합니다.
% Show not-yet-uncovered overlay items transparently.
\setbeamercovered{transparent}

% 노트 페이지의 기본 글자 크기입니다.
% Default font size for note pages.
\setbeamerfont{note page}{size=\footnotesize}

% 발표자 노트를 두 번째 화면 오른쪽에 출력하는 명령입니다.
% Command that shows speaker notes on the right side of a second screen.
\newcommand{\enablenotes}{\setbeameroption{show notes on second screen=right}}

% 발표자 노트를 PDF 출력에서 숨기는 명령입니다.
% Command that hides speaker notes in the PDF output.
\newcommand{\disablenotes}{\setbeameroption{hide notes}}

% 각 슬라이드에 발표자 노트를 쉽게 추가하기 위한 짧은 명령입니다.
% Convenience command for adding speaker notes to each slide.
\newcommand{\slidenote}[1]{\note{\footnotesize #1}}

% 기본값은 청중용 PDF입니다. 노트가 보이지 않습니다.
% Default output is an audience PDF. Notes are hidden.
\disablenotes

% ========== Convenience Macros / 편의 명령 ==========

% 그림, 표, 자료의 출처를 작은 글씨로 표시합니다.
% Display source notes for figures, tables, or materials in tiny text.
\newcommand{\source}[1]{%
  \vspace{0.15em}%
  {\tiny 출처: #1}%
}

% 짧은 직접 인용문과 문헌 인용을 함께 넣을 때 사용합니다.
% Use this for a short quoted phrase followed by a citation.
\newcommand{\qcite}[2]{\enquote{#1} (\parencite{#2})}

% itemize/enumerate의 간격을 줄여 발표 슬라이드에 맞게 만듭니다.
% Tighten itemize/enumerate spacing for presentation slides.
\newcommand{\tightitems}{%
  \setlength{\itemsep}{2pt}%
  \setlength{\parskip}{0pt}%
}


% 테마 색상 선택: blue, green, red 중 하나를 사용합니다.
% Theme color options: choose one of blue, green, or red.
\themecolor{blue}

% 수식 글꼴 선택: rm은 기본 LaTeX 수식 글꼴, sf는 산세리프 계열입니다.
% Equation font option: rm uses the default LaTeX math style; sf uses a sans-serif style.
\equationfont{rm}

% 발표자 노트를 두 번째 화면에 띄우려면 아래 줄의 주석을 해제합니다.
% Uncomment the next line to show speaker notes on a second screen, for example with Pympress.
% \enablenotes

% 새 section이 시작될 때마다 현재 section 중심의 목차 슬라이드를 자동으로 추가합니다.
% Automatically insert a current-section outline slide at the beginning of each section.
\AtBeginSection[]
{
  \begin{frame}{발표 흐름}
  \scriptsize
  % 목차가 길어질 수 있으므로 두 단으로 나누어 표시합니다.
  % Use two columns so long outlines fit on one slide.
  \begin{multicols}{2}
    \tableofcontents[
      currentsection,
      sectionstyle=show/shaded,
      subsectionstyle=show/show/hide
    ]
  \end{multicols}
  \end{frame}
}

% 표지에 들어갈 정보입니다. 실제 값은 sub/preamble_beamer.tex에서 수정합니다.
% Title-page metadata. Edit the actual values in sub/preamble_beamer.tex.
% 대괄호와 texorpdfstring의 두 번째 값은 PDF 메타데이터와 북마크에 사용됩니다.
% The bracketed value and the second texorpdfstring value are used for PDF metadata and bookmarks.
\title[\BeamerShortTitle]{\texorpdfstring{%
  \begin{minipage}{0.92\linewidth}
  \centering
  \BeamerTitleKor\\[0.18cm]
  {\footnotesize\linespread{0.86}\selectfont(\BeamerTitleEng)\par}
  \end{minipage}%
}{\BeamerTitleKor}}
\subtitle{\texorpdfstring{\BeamerSubtitleKor\\{\scriptsize\BeamerSubtitleEng}}{\BeamerSubtitleKor / \BeamerSubtitleEng}}
\author[\BeamerAuthor]{\texorpdfstring{\BeamerAdvisor\\[0.2cm]\BeamerAuthor}{\BeamerAuthor}}
\institute{\BeamerInstitute}
\date{\BeamerDate}

% 실제 PDF 본문이 시작됩니다.
% Start the actual PDF content.
\begin{document}

% 표지 슬라이드입니다.
% Title slide.
\begin{frame}[plain]
  \titlepage
  % 발표자 노트입니다. 기본 설정에서는 PDF에 보이지 않습니다.
  % Speaker note. It is hidden in the PDF by default.
  \slidenote{
    발표 시작 인사와 연구 주제의 핵심 문제를 1분 안에 소개합니다.
    발표 종류, 연구 질문, 방법, 결과 또는 예상 결과가 서로 어떻게 이어지는지 안내합니다.
  }
\end{frame}

% 템플릿을 처음 여는 사용자가 PDF만 읽어도 수정 흐름을 알 수 있도록 넣은 안내 슬라이드입니다.
% This guide slide helps first-time users understand the editing workflow from the generated PDF itself.
\begin{frame}[t]{템플릿 사용 순서}
\scriptsize

\begin{alertblock}{가장 먼저 수정할 곳}
\texttt{sub/preamble\_beamer.tex}에서 아래 값을 바꿉니다.

\texttt{\textbackslash degreeProgramKor}: 석사 또는 박사 \quad
\texttt{\textbackslash presentationStage}: proposal 또는 result

\texttt{\textbackslash BeamerTitleKor}, \texttt{\textbackslash BeamerTitleEng}, \texttt{\textbackslash BeamerAuthor}, \texttt{\textbackslash BeamerAdvisor}
\end{alertblock}

\begin{block}{작성 순서}
\begin{enumerate}
\tightitems
\item 표지 정보와 발표 종류를 정합니다.
\item \texttt{sub/1-Intro\_beamer.tex}부터 \texttt{sub/5-Conclusion\_beamer.tex}까지 예시 문장을 바꿉니다.
\item 그림은 \texttt{images/} 폴더에 넣고 \texttt{\textbackslash includegraphics}로 연결합니다.
\item 참고문헌은 \texttt{sub/references.bib}에 넣습니다.
\item 한국어 문헌은 \texttt{\textbackslash ktextcite}, \texttt{\textbackslash kparencite}; 영어 문헌은 \texttt{\textbackslash textcite}, \texttt{\textbackslash parencite}로 인용합니다.
\item \texttt{latexmk -pdf -lualatex KNUE\_beamer\_main.tex}로 빌드합니다.
\end{enumerate}
\end{block}

\begin{exampleblock}{마무리}
실제 발표 PDF에서는 이 안내 슬라이드와 다음 파일 안내 슬라이드를 삭제하거나 주석 처리합니다.
\end{exampleblock}
\end{frame}

% 파일별 수정 위치를 한 장에 요약합니다.
% Summarize the main editable files in one slide.
\begin{frame}[t]{파일별 수정 위치}
\tiny

\begin{columns}[T]
\begin{column}{0.45\textwidth}
\begin{block}{공통 설정}
\begin{tabular}{@{}p{0.38\linewidth}p{0.54\linewidth}@{}}
\toprule
\textbf{파일} & \textbf{수정 내용} \\
\midrule
\texttt{preamble\_beamer} & 석사/박사, 발표 종류, 제목, 저자, 지도교수, 소속 \\
\texttt{references.bib} & 참고문헌 항목 (한글·영문 공통) \\
\texttt{main.tex} & 발표 순서, 테마 색상, 노트 출력 여부 \\
\bottomrule
\end{tabular}
\end{block}

\begin{alertblock}{변수 선택}
\texttt{\textbackslash degreeProgramKor}, \texttt{\textbackslash presentationStage} 같은 명령 이름은 그대로 두고, 중괄호 안의 값만 바꾸는 방식으로 수정합니다.
\end{alertblock}
\end{column}

\begin{column}{0.45\textwidth}
\begin{block}{본문 슬라이드}
\begin{tabular}{@{}p{0.42\linewidth}p{0.50\linewidth}@{}}
\toprule
\textbf{파일} & \textbf{역할} \\
\midrule
\texttt{1-Intro} & 배경, 필요성, 연구 문제 \\
\texttt{2-Theoretical} & 이론적 배경, 분석틀 \\
\texttt{3-Methods} & 연구 설계, 자료, 분석 방법 \\
\texttt{4-Results} & 예상 결과 또는 연구 결과 \\
\texttt{5-Conclusion} & 요약, 검토 요청 또는 질의응답 \\
\bottomrule
\end{tabular}
\end{block}

\begin{exampleblock}{Tip}
\texttt{proposal} 모드에서는 일정과 예상 산출물을, \texttt{result} 모드에서는 연구 결과와 해석을 중심으로 보여 줍니다.
\end{exampleblock}
\end{column}
\end{columns}
\end{frame}

% 전체 목차 슬라이드입니다.
% Full table-of-contents slide.
\begin{frame}{전체 목차}
  \small
  \tableofcontents[
    sectionstyle=show,
    subsectionstyle=show/show/hide
  ]
  \slidenote{목차는 5개 영역으로 구성합니다. 각 영역은 논문 발표에서 자주 확인하는 질문에 답하도록 배치합니다.}
\end{frame}

% 장별 발표 파일을 불러옵니다. 파일명을 바꾸면 아래 \include 경로도 함께 바꿔야 합니다.
% Include section files. If you rename a file, update the matching \include path below.
% --------------------------------------------------------------------
% 서론 섹션입니다.
% Introduction section.
%
% 이 파일에서는 연구 배경, 연구의 필요성, 연구 목적과 연구 문제를 제시합니다.
% This file presents the research background, rationale, purpose, and questions.
% --------------------------------------------------------------------

% section은 왼쪽 사이드바와 목차에 표시되는 큰 발표 단위입니다.
% A section is a major presentation unit shown in the sidebar and table of contents.
\section{서론}

% subsection은 section 안의 세부 발표 단위입니다.
% A subsection is a smaller unit within a section.
\subsection{연구 배경}

% [t] 옵션은 슬라이드 내용을 위쪽부터 배치합니다.
% The [t] option aligns slide content from the top.
\begin{frame}[t]{연구 배경}

% 슬라이드 본문을 small 크기로 줄여 발표 화면에 안정적으로 맞춥니다.
% Use small text so the slide content fits comfortably.
\small

% block은 일반 설명이나 핵심 개념을 넣는 기본 상자입니다.
% A block is the default box for explanation or key concepts.
\begin{block}{문제 상황}

% itemize는 핵심 내용을 짧은 목록으로 제시할 때 사용합니다.
% Use itemize for concise bullet points.
\begin{itemize}

% tightitems는 목록 간격을 줄이는 사용자 정의 명령입니다.
% tightitems is a helper command that reduces list spacing.
\tightitems

% 아래 항목들은 예시 문장입니다. 자신의 연구 맥락에 맞게 바꿉니다.
% The following bullets are examples. Replace them with your own research context.
\item 연구 분야에서 해결해야 할 교육적, 과학적, 사회적 문제를 제시합니다.
\item 기존 연구나 현장 실천에서 충분히 설명되지 않은 지점을 정리합니다.
\item 연구 대상, 맥락, 핵심 개념을 한 문장으로 연결합니다.
\end{itemize}
\end{block}

% exampleblock은 작성 안내, 장점, 해결 방향처럼 긍정적 안내를 강조할 때 사용합니다.
% Use exampleblock to highlight writing guidance, strengths, or possible solutions.
\begin{exampleblock}{작성 팁}
이 슬라이드는 ``왜 이 연구가 필요한가?''에 답합니다. 배경 설명은 넓게 시작하되 마지막 문장은 자신의 연구 문제로 좁혀 줍니다.
\end{exampleblock}

% slidenote는 발표자 노트입니다. 기본 설정에서는 청중용 PDF에 보이지 않습니다.
% slidenote is a speaker note. It is hidden from the audience PDF by default.
\slidenote{
배경 설명에서는 너무 많은 선행 연구를 나열하기보다, 청중이 연구의 출발점을 빠르게 이해하도록 핵심 문제를 분명하게 말합니다.
}
\end{frame}

% 연구의 필요성에서는 기존 한계와 본 연구의 대응을 나란히 보여 줍니다.
% The rationale slide compares current limitations with the proposed response.
\subsection{연구의 필요성}
\begin{frame}[t]{연구의 필요성}
\small

% columns 환경은 한 슬라이드를 여러 세로 영역으로 나눕니다.
% The columns environment splits one slide into vertical regions.
\begin{columns}[T]

% 첫 번째 열입니다. 전체 본문 너비의 48%를 차지합니다.
% First column. It takes 48% of the text width.
\begin{column}{0.48\textwidth}
\begin{block}{현재 한계}
\begin{itemize}
\tightitems

% 논문 발표에서는 한계를 너무 추상적으로 쓰지 말고 구체적인 공백으로 씁니다.
% In a thesis or dissertation presentation, describe limitations as concrete gaps rather than vague problems.
\item 한계 1: 이론적 공백
\item 한계 2: 방법론적 공백
\item 한계 3: 교육 현장 적용의 어려움
\end{itemize}
\end{block}
\end{column}

% 두 번째 열입니다. alertblock은 문제 해결의 핵심 대응을 강조할 때 적합합니다.
% Second column. alertblock is useful for emphasizing the core response.
\begin{column}{0.48\textwidth}
\begin{alertblock}{본 연구의 대응}
\begin{itemize}
\tightitems

% 각 대응은 왼쪽 한계와 대응되도록 작성하면 설득력이 높아집니다.
% Each response is stronger when it directly answers a limitation in the left column.
\item 새로운 분석 관점 또는 설계 원리 제안
\item 실제 자료, 수업, 사례, 실험 등을 통한 검증
\item 연구 결과의 교육적 활용 가능성 제시
\end{itemize}
\end{alertblock}
\end{column}
\end{columns}

\vspace{0.25cm}

% source 명령은 작은 글씨로 출처나 인용 안내를 표시합니다.
% The source command displays a small source or citation note.
\source{한국어 문헌은 \texttt{\textbackslash kparencite\{kor\_key\}}, 영어 문헌은 \texttt{\textbackslash parencite\{eng\_key\}}로 인용합니다. 예: (\kparencite{kor_seo2009science}; \parencite{eng_nrc2012framework})}

\slidenote{
필요성은 연구자의 관심이 아니라 학문적, 교육적 필요로 설명합니다. 한계와 대응을 나란히 보여 주면 심사자가 연구의 위치를 빨리 파악합니다.
}
\end{frame}

% 연구 목적과 연구 문제는 뒤의 방법, 결과, 결론 전체를 이끄는 중심 슬라이드입니다.
% The purpose and research questions guide the later methods, results, and conclusion.
\subsection{연구 목적과 연구 문제}
\begin{frame}[t]{연구 목적과 연구 문제}
\small

% 연구 목적은 한 문장으로 연구 대상, 처치/분석, 밝힐 내용을 연결합니다.
% The purpose statement links the target, treatment/analysis, and expected explanation in one sentence.
\begin{block}{연구 목적}
본 연구의 목적은 \textbf{[핵심 대상]}을 중심으로 \textbf{[핵심 처치/분석/개발]}을 수행하고, 그 결과 \textbf{[밝히려는 변화 또는 원리]}를 설명하는 데 있다.
\end{block}

% 연구 문제는 번호 목록으로 제시하면 질의응답 때 참조하기 쉽습니다.
% Numbered research questions are easy to reference during Q&A.
\begin{exampleblock}{연구 문제}
\begin{enumerate}
\tightitems

% 연구 문제 수는 실제 연구 범위에 맞게 줄이거나 늘립니다.
% Adjust the number of research questions to match the real research scope.
\item 연구 문제 1: 어떤 특성, 요구, 구조가 확인되는가?
\item 연구 문제 2: 어떤 절차와 원리로 자료, 수업, 모형, 도구를 설계할 수 있는가?
\item 연구 문제 3: 적용 또는 분석 결과 어떤 변화, 효과, 의미가 나타나는가?
\end{enumerate}
\end{exampleblock}

\slidenote{
연구 문제는 뒤의 연구 방법과 1:1로 대응되어야 합니다. 각 연구 문제에 어떤 자료와 분석 방법으로 답할지 자연스럽게 이어지도록 설명합니다.
}
\end{frame}

% --------------------------------------------------------------------
% 이론적 배경 섹션입니다.
% Theoretical background section.
%
% 이 파일에서는 핵심 개념, 선행 연구 흐름, 연구 분석틀을 정리합니다.
% This file organizes key concepts, prior research trends, and the analytical framework.
% --------------------------------------------------------------------

% 발표 목차와 사이드바에 표시될 큰 섹션 이름입니다.
% Major section title shown in the table of contents and sidebar.
\section{이론적 배경}

% 핵심 개념 슬라이드의 하위 섹션입니다.
% Subsection for the key-concepts slide.
\subsection{핵심 개념}

% 핵심 개념은 연구 방법에서 실제 분석 기준으로 다시 사용되어야 합니다.
% Key concepts should later be reused as actual criteria in the methods section.
\begin{frame}[t]{핵심 개념 정리}
\small

% 첫 번째 개념 상자입니다. 연구의 중심 개념을 정의합니다.
% First concept box. Define the central concept of the study.
\begin{block}{개념 1: 연구의 중심 개념}
\begin{itemize}
\tightitems

% 정의는 문헌 정의와 본 연구의 조작적 정의를 구분해 쓰면 좋습니다.
% It is helpful to distinguish literature definitions from your operational definition.
\item 개념의 정의와 범위를 간결하게 제시합니다.
\item 본 연구에서 사용하는 조작적 정의를 함께 적습니다.
\end{itemize}
\end{block}

% 두 번째 개념 상자입니다. 분석틀이나 설계 기준을 넣습니다.
% Second concept box. Use it for the analytical framework or design criteria.
\begin{block}{개념 2: 분석 또는 설계의 기준}
\begin{itemize}
\tightitems

% 여기의 기준은 뒤의 연구 문제, 자료, 분석 방법과 연결되어야 합니다.
% These criteria should connect to the later research questions, data, and analysis methods.
\item 연구 문제를 분석할 때 사용할 이론, 기준, 프레임워크를 설명합니다.
\item 필요하면 표나 개념도로 구성 요소를 정리합니다.
\end{itemize}
\end{block}

\slidenote{
이론적 배경은 연구를 장식하는 문헌 목록이 아니라 분석틀을 세우는 부분입니다. 뒤의 연구 방법에서 이 개념들이 어떻게 쓰이는지 연결해 줍니다.
}
\end{frame}

% 선행 연구는 단순 나열보다 연구 흐름과 남은 공백을 중심으로 구성합니다.
% Prior research should emphasize trends and remaining gaps rather than a simple list.
\subsection{선행 연구}
\begin{frame}[t]{선행 연구 흐름}
\small

% 세 개의 열은 "초기 흐름 - 최근 흐름 - 남은 공백"을 비교하기 위한 구조입니다.
% The three columns compare "early trend - recent trend - remaining gap."
\begin{columns}[T]

% 첫 번째 열: 초기 연구 흐름입니다.
% First column: early research trend.
\begin{column}{0.31\textwidth}
\begin{block}{흐름 1}
\scriptsize
초기 연구가 어떤 문제를 제기했는지 정리합니다.
\end{block}
\end{column}

% 두 번째 열: 최근 연구 흐름입니다.
% Second column: recent research trend.
\begin{column}{0.31\textwidth}
\begin{block}{흐름 2}
\scriptsize
최근 연구가 어떤 방법과 결과를 제시했는지 정리합니다.
\end{block}
\end{column}

% 세 번째 열: 본 연구가 다룰 연구 공백입니다.
% Third column: the research gap that this study will address.
\begin{column}{0.31\textwidth}
\begin{alertblock}{남은 공백}
\scriptsize
아직 충분히 다루지 못한 대상, 맥락, 방법, 자료를 제시합니다.
\end{alertblock}
\end{column}
\end{columns}

\vspace{0.25cm}

% 선행 연구의 흐름도, 개념도, 표를 넣기 전까지 사용할 자리표시자입니다.
% Placeholder for a literature map, conceptual diagram, or summary table.
\placeholderbox{0.72\textwidth}{1.6cm}

\slidenote{
선행 연구는 시간순 나열보다 흐름과 공백 중심으로 제시합니다. 마지막에는 반드시 본 연구가 채울 공백을 말합니다.
}
\end{frame}

% 분석틀은 이론적 배경과 연구 방법을 연결하는 다리 역할을 합니다.
% The analytical framework bridges the theoretical background and research methods.
\subsection{연구 분석틀}
\begin{frame}[t]{연구 분석틀}
\small

% 번호 목록은 분석틀의 구성 요소를 순서 있게 설명할 때 적합합니다.
% A numbered list is useful for presenting framework components in order.
\begin{block}{분석틀 구성}
\begin{enumerate}
\tightitems

% 각 항목은 실제 분석에서 어떤 자료와 연결되는지 설명할 수 있어야 합니다.
% Each item should be explainable in terms of the actual data it will analyze.
\item 이론적 기준: 연구 현상을 해석할 개념 또는 프레임워크
\item 자료 기준: 어떤 자료를 증거로 사용할 것인지에 대한 기준
\item 분석 기준: 코딩, 통계, 비교, 검증 방법의 기준
\end{enumerate}
\end{block}

% 강조 상자는 청중이 이 슬라이드에서 기억해야 할 메시지를 요약합니다.
% The highlighted box summarizes what the audience should remember from this slide.
\begin{exampleblock}{발표에서 강조할 점}
분석틀은 연구 문제와 연구 방법 사이의 다리입니다. 이 슬라이드에서 ``무엇을 어떤 관점으로 볼 것인가''를 분명히 설명합니다.
\end{exampleblock}

\slidenote{
분석틀을 설명할 때는 각 범주가 실제 자료 분석에서 어떻게 판별되는지 예시를 들어 주면 좋습니다.
}
\end{frame}

% --------------------------------------------------------------------
% 연구 방법 섹션입니다.
% Research methods section.
%
% 이 파일에서는 연구 설계, 연구 대상과 자료 수집, 분석 방법을 제시합니다.
% This file presents the research design, participants/data, and analysis methods.
% --------------------------------------------------------------------

% 연구 방법은 연구 문제에 실제로 답하는 절차를 설명하는 섹션입니다.
% The methods section explains the procedures that answer the research questions.
\section{연구 방법}

% 연구 설계 하위 섹션입니다.
% Research design subsection.
\subsection{연구 설계}

% 연구 설계 슬라이드에서는 방법론 이름과 선택 이유를 함께 말합니다.
% In the research design slide, state both the methodology and why it fits.
\begin{frame}[t]{연구 설계}
\small

% 연구 유형을 소개하는 기본 설명 상자입니다.
% Basic explanation box for introducing the research type.
\begin{block}{연구 유형}
\begin{itemize}
\tightitems

% 아래 예시는 자신의 논문 방법론에 맞게 바꿉니다.
% Replace the examples below with the methodology of your own thesis.
\item 예: 설계 기반 연구, 사례 연구, 혼합 연구, 실험 연구, 문헌 연구
\item 연구 유형을 선택한 이유를 연구 문제와 연결해 설명합니다.
\item 설계 기반 연구를 사용할 경우 설계, 실행, 분석, 재설계의 순환 구조를 제시할 수 있습니다(\parencite{eng_dbrc2003design,eng_brown1992design}).
\end{itemize}
\end{block}

% 연구 절차를 시각화하는 간단한 TikZ 흐름도입니다.
% A simple TikZ flow chart visualizing the research procedure.
\begin{center}
\begin{tikzpicture}[node distance=1.0cm, every node/.style={font=\scriptsize}]

% 각 node는 연구 절차의 한 단계를 의미합니다.
% Each node represents one stage of the research procedure.
\node[draw, rounded corners, fill=KNUEPrimaryLight, minimum width=2.1cm] (a) {문제 분석};
\node[draw, rounded corners, fill=KNUEPrimaryLight, minimum width=2.1cm, right=of a] (b) {설계};
\node[draw, rounded corners, fill=KNUEPrimaryLight, minimum width=2.1cm, right=of b] (c) {실행};
\node[draw, rounded corners, fill=KNUEPrimaryLight, minimum width=2.1cm, right=of c] (d) {분석};

% 화살표는 절차의 진행 방향을 나타냅니다.
% Arrows show the direction of the procedure.
\draw[->, thick] (a) -- (b);
\draw[->, thick] (b) -- (c);
\draw[->, thick] (c) -- (d);

% 되돌아가는 화살표는 분석 결과가 재설계로 이어진다는 뜻입니다.
% The returning arrow means analysis results feed into redesign.
\draw[->, thick, bend left=18] (d.north) to (b.north);
\end{tikzpicture}
\end{center}

\slidenote{
방법론 이름만 말하지 말고, 왜 이 연구 문제에 적합한지 설명합니다. 연구 설계 그림은 심사자가 전체 절차를 빠르게 이해하는 데 도움을 줍니다.
}
\end{frame}

% 연구 대상과 자료 수집은 논문 발표에서 구체성을 많이 확인하는 부분입니다.
% Participants and data collection are often checked closely in thesis or dissertation presentations.
\subsection{연구 대상과 자료}
\begin{frame}[t]{연구 대상과 자료 수집}
\small

% 두 열로 나누어 연구 대상과 자료 수집 방법을 구분합니다.
% Use two columns to separate participants/materials from data collection.
\begin{columns}[T]

% 왼쪽 열: 연구 대상 또는 자료 원천입니다.
% Left column: participants or data sources.
\begin{column}{0.48\textwidth}
\begin{block}{연구 대상}
\begin{itemize}
\tightitems

% 참여자 연구라면 표집, 동의, 익명화 계획을 반드시 확인합니다.
% For participant research, check sampling, consent, and anonymization plans.
\item 대상: 학생, 교사, 수업, 문서, 데이터셋 등
\item 표집 기준과 참여자 수
\item 연구 윤리와 동의 절차
\end{itemize}
\end{block}
\end{column}

% 오른쪽 열: 자료를 어떻게 모을지 설명합니다.
% Right column: explain how data will be collected.
\begin{column}{0.48\textwidth}
\begin{block}{자료 수집}
\begin{itemize}
\tightitems

% 연구 문제별로 필요한 자료가 빠지지 않도록 항목을 조정합니다.
% Adjust the items so each research question has the data it needs.
\item 사전/사후 검사 또는 설문
\item 수업 관찰, 산출물, 인터뷰
\item 로그 데이터, 문헌, 공개 데이터
\end{itemize}
\end{block}
\end{column}
\end{columns}

\vspace{0.25cm}

% alertblock은 연구윤리나 누락 위험처럼 주의해야 할 내용을 강조할 때 사용합니다.
% Use alertblock to emphasize cautions such as ethics or missing data.
\begin{alertblock}{주의}
자료 수집 계획은 연구 문제별로 필요한 증거가 빠지지 않도록 구성합니다.
\end{alertblock}

\slidenote{
대상과 자료를 말할 때는 익명화, 동의, 보관 방법 등 윤리적 고려도 함께 언급합니다.
}
\end{frame}

% 분석 방법은 연구 문제, 자료, 분석 절차가 서로 대응되는지 보여 줍니다.
% The analysis slide shows how research questions, data, and methods correspond.
\subsection{분석 방법}
\begin{frame}[t]{분석 방법}
\small

% 표는 연구 문제별 분석 계획을 한눈에 보여 주기 좋습니다.
% A table is useful for showing the analysis plan for each research question.
\begin{block}{연구 문제별 분석 계획}
\begin{tabular}{p{0.28\textwidth}p{0.30\textwidth}p{0.28\textwidth}}

% booktabs의 \toprule, \midrule, \bottomrule은 발표용 표를 깔끔하게 만듭니다.
% booktabs rules make presentation tables cleaner.
\toprule
연구 문제 & 자료 & 분석 방법 \\
\midrule

% 아래 행들은 예시입니다. 실제 연구 문제와 자료명으로 바꿉니다.
% The rows below are examples. Replace them with actual research questions and data.
연구 문제 1 & 설문, 문헌, 사전 자료 & 기술통계, 범주화, 요구 분석 \\
연구 문제 2 & 설계 문서, 전문가 검토 & 내용 타당도, 질적 검토 \\
연구 문제 3 & 산출물, 인터뷰, 검사 결과 & 질적 코딩, 비교 분석, 통계 검정 \\
\bottomrule
\end{tabular}
\end{block}

\slidenote{
이 표는 예시입니다. 실제 연구에서는 각 연구 문제에 답하기 위한 자료와 분석 방법을 구체적으로 바꾸어 작성합니다.
}
\end{frame}

% --------------------------------------------------------------------
% 결과 섹션입니다.
% Results section.
%
% proposal 모드에서는 예상 산출물과 일정을, result 모드에서는 실제 연구 결과와 해석을 제시합니다.
% In proposal mode, present expected outputs and schedule; in result mode, present actual findings and interpretations.
% --------------------------------------------------------------------

% 결과 섹션 제목은 sub/preamble_beamer.tex의 \presentationStage 값에 따라 달라집니다.
% The results-section title changes according to \presentationStage in sub/preamble_beamer.tex.
\section{\BeamerResultsSectionTitle}

% 첫 번째 결과 하위 섹션입니다.
% First results subsection.
\subsection{\BeamerResultsFirstSubsection}

% proposal 모드에서는 예상 산출물을, result 모드에서는 주요 연구 결과를 정리합니다.
% In proposal mode, summarize expected outputs; in result mode, summarize major findings.
\begin{frame}[t]{\BeamerResultsFirstFrame}
\small

% block 안에는 발표 단계에 맞는 결과의 종류를 구체적으로 적습니다.
% In this block, write concrete output or result types for the selected stage.
\begin{block}{\BeamerResultsFirstBlock}
\begin{itemize}
\tightitems

% 아래 예시는 연구 성격에 맞게 삭제하거나 바꿉니다.
% Delete or replace the examples below according to the nature of your study.
\item \KNUEStageText{산출물 1: 분석틀, 수업 모형, 교수-학습자료, 측정 도구 등}{결과 1: 연구 문제 1에 대한 핵심 분석 결과}
\item \KNUEStageText{산출물 2: 적용 결과, 변화 양상, 사례 분석, 설계 원리}{결과 2: 연구 문제 2에 대한 정량/정성 분석 결과}
\item \KNUEStageText{산출물 3: 후속 연구와 현장 적용을 위한 제안}{결과 3: 후속 논의와 현장 적용을 위한 해석}
\end{itemize}
\end{block}

% exampleblock은 발표 단계에 맞는 결과 제시 방식을 안내합니다.
% This exampleblock explains how to present results for the selected stage.
\begin{exampleblock}{\KNUEStageText{결과 제시 방식}{결과 해석 방식}}
\KNUEStageText
{연구계획 발표에서는 확정된 결과보다 \textbf{어떤 자료로 어떤 결과를 얻을 예정인지}와 \textbf{결과가 연구 문제에 어떻게 답할 것인지}를 분명히 보여 줍니다.}
{결과발표에서는 \textbf{무엇을 발견했는지}, \textbf{왜 그렇게 해석할 수 있는지}, \textbf{연구 문제에 어떻게 답하는지}를 분명히 보여 줍니다.}
\end{exampleblock}

\slidenote{
\KNUEStageText
{아직 연구가 끝나지 않았더라도, 예상 산출물의 형태와 판단 기준을 명확히 제시하면 연구의 완성 가능성을 설명할 수 있습니다.}
{결과발표에서는 통계값이나 사례 자체보다 연구 문제와 연결되는 해석을 먼저 말합니다.}
}
\end{frame}

% 두 번째 결과 하위 섹션입니다.
% Second results subsection.
\subsection{\BeamerResultsSecondSubsection}

% proposal 모드에서는 연구 일정을, result 모드에서는 결과 해석을 제시합니다.
% In proposal mode, present the research schedule; in result mode, present result interpretation.
\begin{frame}[t]{\BeamerResultsSecondFrame}
\scriptsize

\KNUEStageText{%
% 표 열 너비는 발표 화면에 맞게 p{...}로 고정했습니다.
% Column widths are fixed with p{...} so the table fits on the slide.
\begin{tabular}{p{0.23\textwidth}p{0.15\textwidth}p{0.48\textwidth}}
\toprule
단계 & 기간 & 주요 내용 \\
\midrule

% 아래 일정은 예시입니다. 실제 학사 일정과 심사 일정을 반영해 수정합니다.
% The schedule below is an example. Revise it according to actual academic and review timelines.
문헌 검토와 분석틀 정리 & 1--2월 & 선행 연구 검토, 연구 문제 확정, 분석 기준 정리 \\
자료 또는 수업 설계 & 3--4월 & 도구 개발, 전문가 검토, 파일럿 점검 \\
자료 수집 & 5--7월 & 수업 실행, 검사, 산출물 수집, 인터뷰 \\
자료 분석 & 8--9월 & 정량 분석, 질적 코딩, 결과 통합 \\
논문 작성과 수정 & 10--12월 & 본문 작성, 지도교수 검토, 최종 수정 \\
\bottomrule
\end{tabular}

\vspace{0.25cm}

% 연구윤리, IRB, 심사 마감처럼 일정에 영향을 주는 요소를 강조합니다.
% Highlight factors that affect the timeline, such as ethics review, IRB, or defense deadlines.
\begin{alertblock}{작성 팁}
실제 학사 일정, 심사 일정, IRB 또는 연구윤리 절차가 있다면 일정표에 함께 반영합니다.
\end{alertblock}

\slidenote{
일정표는 연구가 실행 가능한 계획인지 보여 주는 장치입니다. 각 단계가 연구 문제와 방법에 맞게 이어지는지 설명합니다.
}
}{%
% 결과발표에서는 핵심 결과가 이론, 방법, 선행 연구와 어떻게 연결되는지 설명합니다.
% In result presentations, explain how findings connect to theory, methods, and prior work.
\begin{block}{해석의 초점}
\begin{itemize}
\tightitems
\item 연구 문제별 핵심 결과를 한 문장으로 요약
\item 선행 연구와 일치하거나 다른 지점 설명
\item 방법상의 강점과 제한점이 해석에 미치는 영향 제시
\end{itemize}
\end{block}

\begin{alertblock}{작성 팁}
결과를 나열하기보다, 각 결과가 논문의 주장과 결론을 어떻게 지지하는지 중심으로 설명합니다.
\end{alertblock}

\slidenote{
결과 해석에서는 숫자, 표, 사례를 모두 말하려 하기보다 논문 전체의 주장에 필요한 근거를 선별해 설명합니다.
}
}
\end{frame}

% 세 번째 결과 하위 섹션입니다.
% Third results subsection.
\subsection{\BeamerResultsThirdSubsection}

% proposal 모드에서는 기대 효과를, result 모드에서는 논의와 시사점을 제시합니다.
% In proposal mode, present expected contributions; in result mode, present discussion and implications.
\begin{frame}[t]{\BeamerResultsThirdFrame}
\small

% 두 열 구조로 학문적 기여와 교육적/현장 기여를 나눕니다.
% Use two columns to separate academic and educational/practical contributions.
\begin{columns}[T]

% 왼쪽 열: 학문적 기여입니다.
% Left column: academic contributions.
\begin{column}{0.48\textwidth}
\begin{block}{\KNUEStageText{학문적 기대 효과}{학문적 시사점}}
\begin{itemize}
\tightitems

% 이론, 방법, 후속 연구 측면의 기여를 적습니다.
% Describe contributions to theory, methodology, and future research.
\item \KNUEStageText{이론 또는 분석틀의 확장}{이론 또는 분석틀에 대한 결과의 의미}
\item \KNUEStageText{연구 방법의 적용 가능성 제시}{연구 방법의 타당성과 한계}
\item \KNUEStageText{후속 연구 질문 도출}{후속 연구를 위한 쟁점}
\end{itemize}
\end{block}
\end{column}

% 오른쪽 열: 교육적 또는 현장 적용 기여입니다.
% Right column: educational or practical contributions.
\begin{column}{0.48\textwidth}
\begin{block}{\KNUEStageText{교육적 기대 효과}{교육적 시사점}}
\begin{itemize}
\tightitems

% 수업, 평가, 학습자 지원, 자료 개발 측면의 기여를 적습니다.
% Describe contributions to instruction, assessment, learner support, and material development.
\item \KNUEStageText{수업 설계와 평가 개선}{수업 설계와 평가에 주는 의미}
\item \KNUEStageText{학습자 지원 전략 제안}{학습자 지원 전략에 대한 시사점}
\item \KNUEStageText{현장 적용 자료 제공}{현장 적용 가능성과 주의점}
\end{itemize}
\end{block}
\end{column}
\end{columns}

\slidenote{
\KNUEStageText
{기대 효과는 과장하지 않고 연구 결과가 실제로 기여할 수 있는 범위 안에서 제시합니다.}
{시사점은 결과가 실제로 지지하는 범위 안에서 말하고, 제한점은 결론의 신뢰도를 높이는 방식으로 제시합니다.}
}
\end{frame}

% --------------------------------------------------------------------
% 결론 및 질의 섹션입니다.
% Conclusion and Q&A section.
%
% 이 파일에서는 발표 전체를 요약하고 발표 종류에 맞는 질의 쟁점을 정리합니다.
% This file summarizes the whole presentation and lists Q&A issues for the selected presentation stage.
% --------------------------------------------------------------------

% 마지막 큰 섹션입니다. 발표를 닫고 질의응답으로 넘어갑니다.
% Final major section. It closes the presentation and moves into Q&A.
\section{결론 및 질의}

% 요약 하위 섹션입니다.
% Summary subsection.
\subsection{요약}

% 발표 전체를 한 장으로 다시 연결하는 슬라이드입니다.
% This slide reconnects the whole presentation in one place.
\begin{frame}[t]{\BeamerSummaryFrame}
\small

% 핵심 요약은 연구 문제, 이론, 방법, 결과를 같은 언어로 묶어 줍니다.
% The summary ties research questions, theory, methods, and results together.
\begin{block}{핵심 요약}
\begin{enumerate}
\tightitems

% 아래 문장은 발표자의 실제 연구 내용에 맞게 바꿉니다.
% Replace the following sentences with your actual research content.
\item 연구는 \textbf{[핵심 문제]}에서 출발합니다.
\item 이론적 배경은 \textbf{[핵심 개념/분석틀]}로 정리됩니다.
\item 연구 방법은 \textbf{[대상/자료/분석 방법]}을 통해 연구 문제에 답합니다.
\item \KNUEStageText{예상 결과는 \textbf{[산출물/효과/설계 원리]}로 제시됩니다.}{연구 결과는 \textbf{[핵심 발견/해석/시사점]}으로 정리됩니다.}
\end{enumerate}
\end{block}

% 마무리 문장 예시는 발표 마지막 발화를 준비하기 위한 자리입니다.
% The closing sentence example helps prepare the final spoken message.
\begin{exampleblock}{마무리 문장 예시}
본 연구는 \textbf{[연구 대상]}의 \textbf{[핵심 변화 또는 현상]}를 설명하고, \textbf{[교육적 또는 학문적 기여]}를 제안하고자 합니다.
\end{exampleblock}

\slidenote{
마지막 요약은 처음 제시한 연구 목적과 같은 언어로 닫아 줍니다. 발표 전체가 하나의 논리로 이어졌다는 인상을 주는 것이 중요합니다.
}
\end{frame}

% 질의응답 하위 섹션입니다.
% Q&A subsection.
\subsection{\BeamerQASubsection}

% 심사위원에게 어떤 부분의 피드백이나 질의를 받고 싶은지 명시하는 슬라이드입니다.
% This slide states what kind of feedback or questions you want from the committee.
\begin{frame}[t]{\BeamerQAFrame}
\small

% alertblock으로 질의응답의 핵심 쟁점을 강조합니다.
% Use alertblock to emphasize the key issues for Q&A.
\begin{alertblock}{\BeamerQABlock}
\begin{itemize}
\tightitems

% 실제 발표에서는 지도받고 싶은 쟁점을 구체적인 질문으로 바꿉니다.
% In your actual presentation, turn these into specific questions you want advice on.
\item \KNUEStageText{연구 문제의 범위와 표현이 적절한가?}{연구 결과가 연구 문제에 충분히 답하는가?}
\item \KNUEStageText{연구 방법이 각 연구 문제에 답하기에 충분한가?}{결과 해석이 자료와 분석 방법에 근거하는가?}
\item \KNUEStageText{자료 수집 범위와 분석 기준이 현실적인가?}{논의와 시사점이 결과의 범위를 넘어서지 않는가?}
\item \KNUEStageText{예상 결과와 일정이 학위논문으로 완성 가능한 수준인가?}{최종 논문 수정에서 보완할 핵심 쟁점은 무엇인가?}
\end{itemize}
\end{alertblock}

\vspace{0.2cm}

% 감사 인사와 질의응답 안내 문구입니다.
% Closing thanks and Q&A prompt.
\begin{center}
{\Large 감사합니다.}\\[0.2cm]
{\small 질문과 조언 부탁드립니다.}
\end{center}

\slidenote{
\KNUEStageText
{질의응답에서는 방어적으로 답하기보다, 연구계획을 더 정교하게 만들기 위한 검토 의견으로 받아들이는 태도를 유지합니다.}
{질의응답에서는 결과의 한계와 기여를 균형 있게 설명하고, 최종 논문에 반영할 수정 방향을 정리합니다.}
}
\end{frame}


% 참고문헌 슬라이드입니다. allowframebreaks는 문헌이 많을 때 여러 장으로 자동 분할합니다.
% References slide. allowframebreaks splits long bibliographies across multiple slides.
\begin{frame}[allowframebreaks]{참고문헌}
  \printbibliography
\end{frame}

% 문서 끝입니다.
% End of document.
\end{document}
